% =============================================================================
% Phiên bản conceptual của "Mô hình thực thi ba lớp"
%
% Mục đích: thay thế nội dung hiện tại của \section{Mô hình thực thi ba lớp}
% (chứa các subsubsection chi tiết về Matching Engine, HTLC, Dutch auction
% incentive...) bằng một bản gọn ở mức conceptual. Các chi tiết kỹ thuật
% được chuyển sang mục \S\ref{sec:detailed-components}.
%
% Có thể đặt như \section{...} riêng (giữ vị trí hiện tại của 3.3) hoặc
% đặt như \subsection{...} bên trong "Kiến trúc tổng quan hệ thống" (3.2).
% =============================================================================

\section{Mô hình thực thi ba lớp}
\label{sec:three-layer-model}

Mục 3.2 đã mô tả các \emph{thành phần} của hệ thống và vị trí của chúng
trên sơ đồ off-chain/on-chain. Mục này bổ sung góc nhìn thứ hai —
\emph{chiến lược thực thi} — bằng cách trả lời câu hỏi: \emph{khi một
intent đi vào hệ thống, các thành phần phối hợp với nhau theo thứ tự
nào?}

Cần nhấn mạnh rằng ba lớp dưới đây không phải là ba thành phần kiến
trúc độc lập; chúng là ba \emph{chiến lược thanh toán} áp dụng tuần
tự cho cùng một intent, đều dùng chung tập thành phần đã giới thiệu
ở 3.2.

\subsubsection*{Vấn đề: trade-off giữa chi phí, tốc độ và mức độ
tin cậy}

Các bridge cross-chain hiện hành chia thành hai trường phái chính.
\textbf{Liquidity network} (Across, Stargate, Hop) đặt một bể thanh
khoản trên mỗi chuỗi và thanh toán cho người dùng từ bể đó: tốc độ
cao, trải nghiệm mượt, nhưng phụ thuộc vào LP nên chi phí cao và rủi
ro mất cân bằng pool. \textbf{State validating protocol} (Chainlink
CCIP, LayerZero) dựa trên giả định bảo mật của blockchain nền tảng:
mức độ trust-minimization tốt, nhưng chi phí và độ trễ cao, đặc biệt
giữa các mạng độc lập.

Không chiến lược nào dominate tuyệt đối. Một số intent thuận lợi cho
P2P matching và có thể tự bù trừ với chi phí gần như bằng không;
một số khác xuất hiện trong điều kiện lệch dòng và bắt buộc phải
dùng thanh khoản ngoài. OceanLink chọn tiếp cận \emph{lai}: đặt
nhiều cơ chế thanh toán theo thứ tự ưu tiên và thử lần lượt — thay
vì cố ép mọi intent qua một con đường duy nhất.

\subsubsection*{Ba lớp và thứ tự ưu tiên}

\paragraph{Lớp 1 — Netting Layer.} \emph{Mục tiêu}: thanh toán bằng
cách bù trừ trực tiếp các intent đối ứng, không cần thanh khoản
ngoài. \emph{Cơ chế}: Matching Engine mô hình hoá tập intent đang
chờ thành đồ thị có hướng và tìm các chu trình bù trừ; mỗi chu trình
thoả ngưỡng chấp nhận sẽ được Execution Engine thanh toán nguyên tử
qua HTLC trên các chuỗi liên quan. \emph{Ưu điểm}: chi phí gần bằng
không (chỉ phí gas), trust-minimized vì không cần solver, chất lượng
thực thi tối ưu (không trượt giá). \emph{Nhược điểm}: phụ thuộc vào
sự xuất hiện của intent đối ứng — có thể phải chờ nhiều tick trước
khi tìm được match.

\paragraph{Lớp 2 — Intent Accelerator Layer.} \emph{Mục tiêu}: thanh
toán nhanh phần intent không match được trong lớp 1, bằng cách mời
các solver bên ngoài cạnh tranh quote. \emph{Cơ chế}: phần dư của
intent được broadcast đến các solver đã đăng ký; mỗi solver đề xuất
một mức giá riêng, người gửi chấp nhận quote tốt nhất. Một biến thể
quan trọng là \emph{incentive Dutch auction}: phí thưởng cho solver
bắt đầu bằng không và tăng dần theo thời gian. Điều này tạo trò chơi
chiến lược — solver phải cân nhắc giữa thực thi sớm với phần thưởng
thấp hay chờ phần thưởng cao hơn nhưng đối mặt với rủi ro mất lượt.
\emph{Ưu điểm}: tốc độ thực thi gần như tức thời, tăng đáng kể tỉ
lệ hoàn tất. \emph{Nhược điểm}: chi phí cao hơn lớp 1 (incentive
trả cho solver), tiềm năng các hành vi chiến lược như selective
fill.

\paragraph{Lớp 3 — Fallback Bridge Aggregator.} \emph{Mục tiêu}: bảo
đảm completion cho intent vẫn còn dư sau hai lớp trên, bằng cách
định tuyến qua một bridge bên thứ ba có thanh khoản thực sự (Across,
LayerZero, Chainlink CCIP). \emph{Cơ chế}: aggregator gọi
\texttt{quote} song song trên các adapter bridge, chọn quote tốt
nhất theo hàm chi phí (\,$\text{phí} + \alpha \cdot \text{ETA}$\,)
và phát giao dịch qua bridge tương ứng. \emph{Ưu điểm}: hoàn tất
gần như chắc chắn, không phụ thuộc dòng intent nội bộ.
\emph{Nhược điểm}: chi phí cao nhất trong ba lớp, độ trễ phụ thuộc
bridge bên ngoài, và execution path dài làm việc theo dõi phức tạp.

\subsubsection*{Vòng đời một intent qua ba lớp}

Mỗi intent đi vào hệ thống đều khởi đầu ở Lớp 1 và \emph{leo dần}
xuống các lớp dưới khi cần:

\begin{enumerate}
  \item Intent được \texttt{validateIntent} ở API Gateway và đưa vào
    OrderStore với trạng thái \texttt{QUEUED}.
  \item Mỗi tick scheduler (mặc định 5\,s) Matching Engine cố tìm
    cycle trên toàn bộ batch intent đang sống. Trong cửa sổ thời gian
    $T_{\text{net}}$:
    \begin{itemize}
      \item Nếu tìm được cycle khớp toàn bộ: trạng thái chuyển
        \texttt{MATCHED}, Execution Engine settle nguyên tử qua HTLC,
        intent kết thúc ở Lớp 1.
      \item Nếu khớp một phần: phần khớp được settle, phần dư trở
        thành \texttt{PARTIAL} và tiếp tục chờ.
    \end{itemize}
  \item Khi cửa sổ $T_{\text{net}}$ trôi qua mà phần dư vẫn còn,
    intent (hoặc phần dư của nó) được chuyển sang Lớp 2: broadcast
    cho solver, chạy Dutch auction. Nếu một solver thực thi, phần
    đó settle; nếu không, sang Lớp 3.
  \item Lớp 3 chọn bridge tốt nhất, phát giao dịch và theo dõi đến
    khi hoàn tất hoặc thất bại. Đây là điểm cuối — không có lớp 4.
\end{enumerate}

Trật tự "leo dần" này đảm bảo phần lớn intent hưởng chi phí thấp
nhất có thể (Lớp 1), một số ít chấp nhận trade-off cho tốc độ (Lớp
2), và \emph{không có intent nào bị treo vô thời hạn} (Lớp 3 luôn
sẵn sàng làm điểm thoát).

\subsubsection*{Ánh xạ vào các thành phần}

Bảng dưới đây liên kết mỗi lớp với các thành phần được kích hoạt:

\begin{itemize}
  \item \textbf{Lớp 1} dùng API Gateway, OrderStore, Matching Engine
    (MatchScheduler + MatchingService + Solver), Execution Engine và
    Smart Contract Layer.
  \item \textbf{Lớp 2} dùng cùng API Gateway và OrderStore, nhưng
    thay Matching Engine bằng một bộ phát quote tới các solver bên
    ngoài; Execution Engine và HTLC vẫn là kênh thanh toán cuối.
  \item \textbf{Lớp 3} dùng Fallback Bridge Aggregator để chọn và
    gọi adapter bridge phù hợp; bỏ qua HTLC và đi qua kênh thanh
    toán riêng của bridge bên thứ ba.
\end{itemize}

Như vậy, các thành phần kiến trúc đã trình bày ở 3.2 \emph{phục vụ
song song} cho cả ba lớp; ba lớp khác nhau ở \emph{thứ tự huy động}
và \emph{cơ chế tìm đối tác}, không ở mặt cấu trúc phần mềm. Chi
tiết mỗi thành phần được trình bày ở mục
\S\ref{sec:detailed-components}.

\subsubsection*{Phạm vi triển khai trong luận văn}

Trong scope của luận văn:

\begin{itemize}
  \item \textbf{Lớp 1} đã được cài đặt đầy đủ: matching engine, HTLC
    multi-fill, orchestrator. Đây là phần tập
    trung của các mục thiết kế chi tiết và phần đo benchmark trong
    chương Đánh giá.
  \item \textbf{Lớp 2} được trình bày ở mức thiết kế: interface,
    Dutch auction incentive, lưu đồ tương tác. Cài đặt cụ thể
    nằm ngoài scope.
  \item \textbf{Lớp 3} được trình bày ở mức thiết kế: interface
    \texttt{BridgeAdapter}, danh sách bridge ứng cử và lưu đồ
    quote–send. Cài đặt cụ thể nằm ngoài scope.
\end{itemize}

Phần Future Work (chương cuối) liệt kê các hạng mục cần hoàn thiện
để Lớp 2 và Lớp 3 đi vào hoạt động đầy đủ.
